\section{Files this design touches}\label{sec:files}

This section enumerates the concrete edits the refactor makes to the working tree, mapping
each file to its change and then bounding the change's reach. It serves two purposes: it is
the implementer's manifest, and it makes the surface area of the change auditable up front.
Every line number below was verified against the current tree (branch \texttt{PartPlacement},
namespace \code{model::part}); the design specification's own citations were pinned to an
older commit and several have since drifted, so the table reflects the live source, not the
spec. The defining property of the design is that it is \emph{narrow}: a single in-slot
storage swap, one new header, and a fixed set of consumer rewrites. There is no change to the
force path, the integrator, or the state vector.

\Cref{tab:files} is the file-by-file change map. The intent of presenting it this way is to
make explicit that almost every edit is either (a) adding a closed-form geometry virtual to a
part that already stores the backing radii, or (b) swapping a consumer from the legacy
CM-to-CM offset to the resolver's geometric \code{Pose}. The single genuinely new
artifact is \srcf{model/parts/Placement.h} (\cref{sec:datamodel}); everything else is an
amendment to existing code.

\begin{table}[htbp]
\caption{Files this design touches, with the change applied to each. Line numbers are pinned
to the current working tree, not the design spec. \emph{New} marks the one wholly new file;
all other entries amend existing code.}\label{tab:files}
\small
\begin{tabularx}{\linewidth}{@{}>{\raggedright\arraybackslash}p{0.27\linewidth}>{\raggedright\arraybackslash}X@{}}
\toprule
\textbf{File} & \textbf{Change} \\
\midrule
\src{model/parts/Part.h}{326}
  & Storage member \code{childParts} changes type from
    \code{std::vector<std::tuple<\allowbreak std::shared\_ptr<Part>,\allowbreak Vector3>>} to
    \code{std::vector<std::pair<\allowbreak std::shared\_ptr<Part>,\allowbreak StationLink>>} --- a single
    in-slot swap (\cref{sec:datamodel}). Adds the geometry virtuals \code{stationAt},
    \code{radiusOuterAt}, \code{radiusInnerAt}, \code{axialLength}, \code{isSolid}, and
    \code{cmStationLocal} to the base class (none exist today; the only base geometry virtual
    is \code{getReferenceArea}, \src{model/parts/Part.h}{176}). Promotes the new
    \code{StationLink} \code{addChildPart} overload (\src{model/parts/Part.h}{239}) and
    updates the \code{getChildParts} return type (\src{model/parts/Part.h}{212}). Adds a
    \code{placementDirty} flag and a \code{resolvedCache} member. \\
\addlinespace
\srcf{model/parts/Placement.h}
  & \emph{New.} Houses the value types and free functions: \code{Station}, \code{StationLink},
    \seat{SeatKind}, \code{Pose}, \code{Placed}, \code{resolvePlacements},
    \code{sweepOverlaps}, \code{OverlapDiagnostic}, and \code{SolveResult}
    (\cref{sec:datamodel,sec:resolver,sec:diagnostics}). \\
\addlinespace
\srcf{model/parts/Part.cpp}
  & \code{computeCompositeAt} (\src{model/parts/Part.cpp}{169}) is rewritten to consume the
    resolver's \code{Pose} and derive CM through \code{cmStationLocal}, preserving the
    two-pass mass-weighted CM and parallel-axis structure (\cref{sec:resolver}).
    \code{accumulateAeroAt} (\src{model/parts/Part.cpp}{214}) threads each child's resolved
    \code{Pose.origin.z()} instead of \code{axialStation + pos.z()}
    (\src{model/parts/Part.cpp}{224}). Adds \code{ensurePlacementCache}. Every structured
    binding over \code{childParts} is updated from the tuple to the pair (see below). \\
\addlinespace
\srcf{model/parts/ConicalNoseCone.cpp}
  & Adds \code{radiusOuterAt} (the linear taper $\text{baseRadius}\cdot(1-z/L)$), the
    \code{cmStationLocal} frame adapter of \cref{sec:datamodel} (the cone is the only part
    that overrides it), and \code{isSolid}. The internal cone CM math in \code{coneCmOffset}
    (\src{model/parts/ConicalNoseCone.cpp}{49}) is \emph{not} touched; only the adapter that
    reconciles its frame is added. \\
\addlinespace
\srcf{model/parts/BodyTube.h}, \srcf{model/parts/BodyTube.cpp},
\srcf{model/parts/FinSet.h}, \srcf{model/parts/FinSet.cpp},
\srcf{model/parts/HollowSphere.h}, \srcf{model/parts/HollowSphere.cpp}
  & Each gains \code{radiusOuterAt} / \code{radiusInnerAt} and \code{isSolid} backed by its
    existing wall radii (\code{BodyTube}: \src{model/parts/BodyTube.h}{49};
    \code{HollowSphere}: \src{model/parts/HollowSphere.h}{52}). \code{FinSet} reports its
    \emph{body-disc} radius only (\code{getBodyRadius}, \src{model/parts/FinSet.h}{67}) for
    axial occupancy; sectored fin interference is scoped out (\cref{sec:diagnostics}). \\
\addlinespace
\srcf{visualizer/RocketMesh.cpp}
  & \code{walk} (\src{visualizer/RocketMesh.cpp}{427}) consumes the resolver's \code{Pose}
    instead of accumulating \code{station + offset} (\src{visualizer/RocketMesh.cpp}{472}).
    Because the convention is \zfwd{} --- already the visualizer's own frame
    (\src{visualizer/RocketMesh.cpp}{156}) --- this is a pure consumer swap with \textbf{no
    axis flip} (\cref{sec:philosophy}). \\
\addlinespace
\srcf{model/DesignSerializer.cpp}
  & \code{writeOffset} (\src{model/DesignSerializer.cpp}{53}) becomes \code{writeLink},
    emitting \code{seat}/\code{parentStation}/\code{childStation}/\code{gap} instead of an
    \code{x}/\code{y}/\code{z} offset. \code{buildPart}
    (\src{model/DesignSerializer.cpp}{117}) reads a \code{<link>} element. The version
    string bumps from \texttt{0.1} (\src{model/DesignSerializer.cpp}{156}) to \texttt{0.2};
    the loader's \code{major == "0"} tolerance (\src{model/DesignSerializer.cpp}{186})
    already accepts it, so no gate change is needed (\cref{sec:serialization}). \\
\bottomrule
\end{tabularx}
\end{table}

\subsection*{The tuple-to-pair destructure sites in \texttt{Part.cpp}}

Changing the element type of \code{childParts} mechanically invalidates every structured
binding that walks it. These are local, internal sites --- all within
\srcf{model/parts/Part.cpp} --- and each must be re-spelled to bind a
\cppinline|pair<shared_ptr<Part>, StationLink>| (and, where it matters, to read the resolved
geometric station rather than the stored offset). The complete, verified set is: \code{clone}
(\src{model/parts/Part.cpp}{117}, including the \code{emplace\_back} that reconstructs the
element at \src{model/parts/Part.cpp}{121}), \code{getCompositeMass}
(\src{model/parts/Part.cpp}{131}), \code{computeCompositeAt} pass~1
(\src{model/parts/Part.cpp}{179}), \code{accumulateAeroAt}
(\src{model/parts/Part.cpp}{221}), \code{maxFrontalReferenceArea}
(\src{model/parts/Part.cpp}{231}), \code{findById} (\src{model/parts/Part.cpp}{244}), and
both loops in \code{removeChildById} --- the \cppinline|std::get<0>| direct-child scan
(\src{model/parts/Part.cpp}{261}) and the recursive descent
(\src{model/parts/Part.cpp}{272}).

\begin{keyinsight}[Two sites the spec's enumeration mis-states]
The design spec lists the destructure sites as lines \texttt{179, 197, 221, 231, 244, 261,
272}. Verification against the live tree corrects this in two ways. First, line~197 is
\emph{not} a \code{childParts} site: it destructures the local
\cppinline|std::vector<std::pair<Vector3, CompositeProperties>>| accumulator in
\code{computeCompositeAt} pass~2, which is unaffected by the storage swap. Second, the spec
omits two real \code{childParts} sites entirely --- \code{clone}
(\src{model/parts/Part.cpp}{117}) and \code{getCompositeMass}
(\src{model/parts/Part.cpp}{131}). Both bind the element shape and so both must be migrated.
An implementer working from the spec's list alone would leave \code{clone} and
\code{getCompositeMass} broken.
\end{keyinsight}

\subsection*{Blast radius: the public \texttt{getChildParts()} return type}

The destructure sites above are private to \code{Part}. The change with reach beyond the class
is the public accessor: \code{getChildParts} (\src{model/parts/Part.h}{209}) currently returns
\cppinline|const std::vector<std::tuple<std::shared_ptr<Part>, Vector3>>&|, and the refactor
changes that to the \code{pair<ptr, StationLink>} vector. Every external consumer that iterates
the returned view therefore breaks at compile time and must be migrated --- a
\cppinline|std::get<0>| $\rightarrow$ \code{.first} (and \cppinline|std::get<1>| $\rightarrow$
\code{.second}) rename, or the equivalent structured-binding rewrite. This is mechanical and
enumerated, not hidden.

\begin{designrule}[The shim does not cover the accessor]
The \cppinline|[[deprecated]]| transitional \code{addChildPart(child, Vector3)} overload
(\cref{sec:authoring}) keeps legacy \emph{setter} call sites compiling through the migration.
It does \textbf{not} shield the \emph{accessor}: \code{getChildParts}'s element type is part of
the public contract, and there is no way to return both shapes from one signature. Accessor
consumers must therefore be migrated explicitly, in lockstep with the storage swap, rather than
deferred behind the shim.
\end{designrule}

The confirmed external call sites are the visualizer's tree walk
(\src{visualizer/RocketMesh.cpp}{466}, which today reaches the element through
\cppinline|std::get<0>| / \cppinline|std::get<1>|) and the serializer's child-writing loop
(\src{model/DesignSerializer.cpp}{73}); both files already appear in \cref{tab:files} for their
primary changes, so the accessor migration folds into work those files receive regardless.
Beyond these, the CLI and GUI tree-walkers and any tests that iterate \code{getChildParts}
update the same way.

\subsection*{A new base virtual: \texttt{getLength()}}

One amendment in \cref{tab:files} is easy to overlook because it is implied rather than stated
as a line edit: \code{getLength()} is \emph{promoted to a base \code{Part} virtual}. Today
there is no \code{getLength()} on \code{Part} at all. It exists only as a non-virtual accessor
on the concrete types that have an axial length --- \code{ConicalNoseCone}
(\src{model/parts/ConicalNoseCone.h}{53}) and \code{BodyTube}
(\src{model/parts/BodyTube.h}{51}) --- while \code{FinSet} and \code{HollowSphere} have none.
The base \code{stationAt} and the default \code{axialLength()} both need a part-level length to
resolve fractional stations, so the refactor lifts \code{getLength()} onto the base class.

\code{HollowSphere} consequently gains a \code{getLength()} override: a sphere has no length
member (\src{model/parts/HollowSphere.h}{52} exposes only radii), so its axial extent must be
reported as its diameter, $2\cdot\text{outerRadius}$. Without this override the inherited
default \code{axialLength()} would have nothing to read for a sphere. This is the smallest of
the file changes, but it is the one that makes the fractional-station model
(\cref{sec:datamodel,sec:resolver}) well defined for every concrete part type rather than only
the two that happened to carry a length already.
